%! Regression Lines and Residuals — Unit 2, Lesson 2.3 — Warm-Up
% ONE PAGE, blank and key. Three items at 12pt. Do not add a fourth.
% GRADUAL RELEASE: the warm-up activates prior knowledge before the guided notes.
% Item 1 is the spiral to Lesson 2.2 (explanatory variable, direction, strength of r) on the
% pizza-delivery data the notes open on; item 2 is plugging into a line — exactly the move
% notes row 1 names "prediction"; item 3 is "above or below the line, and by how much" — the
% residual, in algebra language, before the notes name it and fix its sign.
\documentclass[12pt]{article}
\usepackage{apstats-article}
\usepackage{apstats-boxes}

\begin{document}

\pageheader{Unit 2, Lesson 2.3}{Warm-Up}

\vspace{0.28in}

\begin{enumerate}[leftmargin=1.9em, itemsep=1.0in]

  \item \textbf{From last lesson.} A pizza shop plotted delivery distance (miles) against
        delivery time (minutes) for 17 deliveries. The points rise from left to right in a
        tight, straight band.

        \vspace{12pt}
        Explanatory variable: \blank{3.4cm} \qquad Direction: \blank{2.4cm}

        \vspace{16pt}
        Is $r$ closer to $0.3$ or to $0.95$? \blank{1.6cm} \qquad Form: \blank{2.4cm}

  \item Use the line $y = 8 + 3x$.

        \vspace{12pt}
        When $x = 4$, \quad $y =$ \blank{3.6cm} \qquad When $x = 7$, \quad $y =$ \blank{3.6cm}

  \item The point $(5,\ 26)$ is plotted beside the line $y = 8 + 3x$.

        \vspace{12pt}
        At $x = 5$ the line is at height \blank{3.4cm}

        \vspace{16pt}
        Is the point above or below the line? \blank{2.0cm} \quad By how much? \blank{1.6cm}

        \vspace{16pt}
        Now the point $(4,\ 17)$: above or below? \blank{2.0cm} \quad By how much? \blank{1.6cm}

        \vspace{16pt}
        Which point's actual $y$ came out \emph{bigger} than the line predicted?\\[18pt]
        \writeline

\end{enumerate}

\end{document}
